\section{Del RL multitarea al metaaprendizaje}

\subsection{Repaso del RL multitarea}

En la clase anterior se trató el RL multitarea: se aumenta el estado a $(s, z)$ (donde $z$ es la descripción de la tarea) y se entrena una política condicionada $\pi(a|s,z)$ para resolver varias tareas. Sin embargo, el RL multitarea tiene una limitación importante: supone que \textbf{todas las tareas se conocen desde el entrenamiento}.

\begin{figure}[H]
\centering
\includegraphics[width=0.9\textwidth]{slides-images/slide-03.jpg}
\caption{Repaso del RL multitarea: la descripción de la tarea se incorpora al estado y todo se unifica en un joint MDP}
\end{figure}

\begin{knowledgebox}{Limitaciones del RL multitarea}
El RL multitarea cubre todas las tareas con una sola política compartida. Si en la prueba aparece una tarea nueva muy distinta de las tareas de entrenamiento, el desempeño puede caer drásticamente. Además, a medida que aumenta el número de tareas, se vuelve cada vez más difícil aprender una política que funcione bien en todas (el problema de interferencia entre tareas).
\end{knowledgebox}

\begin{figure}[H]
\centering
\includegraphics[width=0.9\textwidth]{slides-images/slide-04.jpg}
\caption{El goal-conditioned learning es un caso particular del RL multitarea}
\end{figure}

\subsection{Idea central del metaaprendizaje}

\begin{importantbox}{Definición de Meta-RL}
El objetivo del metaaprendizaje por refuerzo (Meta-RL) no es aprender «cómo resolver una tarea concreta», sino \textbf{aprender «cómo aprender rápidamente tareas nuevas»}.

Formalmente: dada una distribución de tareas $p(\mathcal{T})$, en la etapa de entrenamiento el Meta-RL usa muchas tareas muestreadas de $p(\mathcal{T})$ para aprender un «algoritmo de aprendizaje» o un «mecanismo de adaptación rápida». En la prueba, ante una tarea nueva $\mathcal{T}_{\text{new}} \sim p(\mathcal{T})$, basta con pocos datos para adaptarse.
\end{importantbox}

\subsection{Por qué se necesita el Meta-RL}

El ponente ilustra el problema con un ejemplo de la vida cotidiana: aprender a usar una máquina de espresso nueva. Si esta tarea se le asigna a un agent de RL entrenado desde cero, podría necesitar millones de intentos; en cambio, a una persona le bastan un poco de experiencia relacionada y algunas instrucciones para dominarla en unos minutos. La diferencia no está en que el ser humano sepa «optimizar la recompensa» mejor, sino en que el ser humano \textbf{nunca empieza desde cero}.

\begin{figure}[H]
\centering
\includegraphics[width=0.9\textwidth]{slides-images/slide-06.jpg}
\caption{Motivación del metaaprendizaje por refuerzo: ante una tarea nueva, las personas aprovechan su experiencia previa para adaptarse con rapidez}
\end{figure}

\begin{knowledgebox}{La carencia que el Meta-RL intenta cubrir}
El RL convencional supone que cada tarea nueva exige recolectar de nuevo una gran cantidad de datos de interacción; el Meta-RL, en cambio, intenta codificar en el sistema de aprendizaje el hecho de «haber resuelto antes tareas similares», para que en una tarea nueva alcance más rápido un régimen de comportamiento eficaz.
\end{knowledgebox}

\subsection{Del transfer learning al meta-learning}

En esta clase también se comparan varios paradigmas de transferencia comunes:
\begin{enumerate}
    \item \textbf{Forward transfer}: primero se resuelve la source task y después se hace fine-tuning hacia la target task.
    \item \textbf{Multi-task transfer}: se entrena simultáneamente en varias tareas y se usa un task descriptor para permitir la generalización sin ejemplos.
    \item \textbf{Meta-learning}: desde la etapa de entrenamiento se optimiza explícitamente «cómo adaptarse rápidamente a tareas nuevas».
\end{enumerate}

\begin{figure}[H]
\centering
\includegraphics[width=0.9\textwidth]{slides-images/slide-07.jpg}
\caption{Diferencias en el planteamiento del problema entre transfer, multi-task y meta-learning}
\end{figure}

\begin{warningbox}{Límites de generalización del Meta-RL}
El Meta-RL no garantiza la adaptación a cualquier tarea nueva. Depende de un supuesto fundamental: las tareas nuevas que aparecen en el meta-test siguen proviniendo de una distribución de tareas cercana a la del meta-train. Si las tareas de prueba quedan por completo fuera de la distribución de entrenamiento, tanto la velocidad de adaptación como el desempeño final pueden disminuir de forma significativa.
\end{warningbox}

\subsection{La perspectiva del aprendizaje few-shot}

El curso también plantea una analogía con la few-shot image classification y el in-context learning de los LLM. El punto en común es que el modelo no recibe solo la entrada $x$, sino también un pequeño conjunto de ejemplos $D_{\text{train}}$, y a partir de él infiere a qué tipo de tarea pertenece el problema actual y cómo debe responder.

\begin{figure}[H]
\centering
\includegraphics[width=0.9\textwidth]{slides-images/slide-09.jpg}
\caption{Abstracción del few-shot learning: la función de predicción depende a la vez de la entrada y de unos pocos ejemplos de entrenamiento}
\end{figure}

\subsection{Resumen de la sección}

El Meta-RL eleva el nivel del aprendizaje de «aprender una política» a «aprender a aprender políticas», y es adecuado para escenarios que exigen adaptarse rápidamente a tareas nuevas.
